%==============================================================================
%  CAPÍTULO IV
%==============================================================================
\chapter{Procedimiento concursal de reorganización judicial ordinario}
\label{cap:reorganizacion-ordinaria}

\fichaCapitulo{Salvataje de la mediana y gran empresa.}{Tramitación de la \pfc{}
ordinaria, quórums de aprobación, financiamiento durante la protección y régimen de
impugnación del acuerdo.}

%==============================================================================
\bloqueTeorico
%==============================================================================

\subsection{Conservación de la empresa viable frente a la liquidación forzosa}

La reorganización judicial descansa sobre una premisa económica verificable: existe una
categoría de empresas cuyo valor como unidad en funcionamiento excede el valor de
liquidación de sus activos considerados aisladamente. La diferencia entre ambos valores
---lo que la literatura denomina valor de empresa en marcha--- se compone de elementos que
no sobreviven al desmembramiento: la organización del trabajo, la cartera de clientes, los
contratos vigentes, el conocimiento acumulado y la reputación comercial.

Cuando una empresa de esa categoría entra en crisis financiera ---esto es, cuando no puede
servir su pasivo en las condiciones pactadas, pero su operación sigue generando margen---,
la liquidación destruye valor que nadie captura. Los acreedores reciben menos de lo que
habrían obtenido reestructurando; los trabajadores pierden empleos productivos; los
proveedores pierden un cliente solvente en potencia.

La reorganización existe para evitar esa destrucción. No es un beneficio concedido al
deudor: es un mecanismo de maximización del valor disponible para los acreedores, que
opera precisamente porque son ellos quienes deciden si el plan les conviene más que la
liquidación.

\begin{foco}[El criterio que ordena todo el capítulo]
La pregunta que debe responder cualquier propuesta de acuerdo es una sola: ¿ofrece a los
acreedores más de lo que obtendrían en un escenario de liquidación? Si la respuesta es
afirmativa y está documentada, el acuerdo tiene fundamento económico. Si no lo es, ningún
argumento jurídico lo salvará en la junta.
\end{foco}

\subsection{Naturaleza de la Protección Financiera Concursal}

La \pfc{} es una interferencia deliberada del legislador en el derecho de los acreedores a
perseguir individualmente al deudor. Durante su vigencia, quien tiene título ejecutivo no
puede ejecutar; quien tiene derecho a poner término a un contrato por incumplimiento no
puede hacerlo por esa causa; quien podría solicitar la liquidación no puede pedirla.

La justificación de una interferencia tan intensa es funcional y transitoria. Sin ella la
reorganización sería imposible: bastaría que un acreedor individual ejecutara un activo
esencial para que la empresa dejara de operar, y con ello desapareciera el valor que se
intentaba preservar. La protección crea el espacio temporal en que la negociación
colectiva puede desarrollarse.

De su carácter instrumental derivan dos límites que recorren todo el régimen: la protección
es \emph{temporal}, con un techo legal infranqueable, y es \emph{parcial}, en cuanto no
alcanza a todos los acreedores por igual.

\subsection{Ámbito subjetivo: sólo la Empresa Deudora}

El \art{54} restringe el procedimiento a la \empresadeudora{}, que en el Capítulo III de la
ley se denomina indistintamente Empresa Deudora o Deudor. La \personadeudora{} carece de
acceso a esta vía: su equivalente funcional es la renegociación administrativa (capítulo
\ref{cap:renegociacion}).

El procedimiento se inicia mediante solicitud presentada por la propia empresa ante el
tribunal correspondiente a su domicilio. No existe reorganización forzosa: la iniciativa es
exclusivamente del deudor, lo que constituye una diferencia estructural con la liquidación,
que admite la iniciativa del acreedor.

%==============================================================================
\bloqueNormativo
%==============================================================================

\subsection{Inicio del procedimiento y antecedentes de la solicitud}

\subsubsection{La solicitud y su modelo obligatorio}

El \art{54} encomienda a la \superir{} regular mediante norma de carácter general el modelo
de solicitud, que debe estar disponible para los interesados. Ese modelo es la \ncg{4}
\citeyearpar{ncg-4}, la norma vigente más antigua del sistema.

\begin{practica}[Usar el modelo no es opcional]
La solicitud que se aparta del modelo aprobado expone al deudor a una objeción que consume
días de un procedimiento cuyo plazo de protección ya está corriendo. Antes de redactar debe
descargarse el modelo vigente y verificarse su fecha.
\end{practica}

\subsubsection{Los antecedentes del artículo 56}

Aceptada la nominación por el veedor titular y suplente, la Superintendencia remite al
tribunal el Certificado de Nominación. Paralelamente, el deudor debe acompañar, mediante
declaración jurada simple firmada, los antecedentes que enumera el \art{56}. Los dos
primeros concentran la mayor densidad de información y de riesgo:

\begin{enumerate}
  \item \textbf{Relación de todos sus bienes}, con expresión de su avalúo comercial, del
  lugar en que se encuentren y de los gravámenes que los afecten. Debe señalarse, además,
  \destacar{cuáles de esos bienes tienen la calidad de esenciales para el giro}.

  \item \textbf{Relación de los bienes de terceros constituidos en garantía} en favor del
  deudor, con la misma indicación de esencialidad.
\end{enumerate}

\begin{foco}[La calificación de esencialidad es una decisión estratégica]
Declarar un bien como esencial para el giro lo sustrae de la libre disposición y refuerza el
argumento de continuidad operativa. Declararlo no esencial abre la posibilidad de enajenarlo
para generar liquidez durante la protección. La calificación se hace una sola vez, bajo
declaración jurada, y condiciona todo el procedimiento posterior. Debe decidirse con el
flujo proyectado a la vista, no como un trámite de llenado.
\end{foco}

\begin{alerta}[El avalúo comercial y su contraste]
El \art{56} exige avalúo \emph{comercial}, no valor libro ni valor de adquisición. Un avalúo
inflado deteriora la credibilidad de toda la declaración cuando el veedor lo contraste, y en
caso de liquidación posterior puede constituir el antecedente de la causal del \art{169 A}
\Nro 1 (capítulo \ref{cap:mala-fe}).
\end{alerta}

\subsection{La Resolución de Reorganización}

\subsubsection{Plazo y contenido}

Conforme al \art{57}, dentro del quinto día de efectuada la presentación el tribunal dicta
la Resolución de Reorganización, en la que designa a los veedores titular y suplente
nominados según el \art{22} y dispone los efectos de la protección.

\subsubsection{Efectos de la protección}

El \art{57} \Nro 1 enumera los efectos:

\begin{description}[leftmargin=0pt,style=nextline,font=\sffamily\bfseries]
  \item[Letra a) — Inmunidad frente a la ejecución] No puede declararse ni iniciarse un
  procedimiento concursal de liquidación en contra del deudor, ni iniciarse en su contra
  juicios ejecutivos, ejecuciones de cualquier clase o restituciones en juicios de
  arrendamiento.

  \item[Letra b) — Suspensión] Se suspende la tramitación de los procedimientos señalados y
  los plazos de prescripción extintiva.

  \item[Letra c) — Vigencia de los contratos] Todos los contratos suscritos por el deudor
  mantienen su vigencia y sus condiciones de pago, y no pueden terminarse anticipadamente
  por la sola circunstancia del procedimiento.
\end{description}

\begin{alerta}[La excepción laboral: el límite más importante de la protección]
La inmunidad de la letra a) \destacar{no alcanza a los juicios laborales sobre obligaciones
que gocen de preferencia de primera clase}. Respecto de ellos se suspende únicamente la
ejecución y realización de bienes del deudor: el juicio continúa tramitándose.

La excepción tiene una contraexcepción antifraude. No opera cuando se trata de juicios
laborales de esa clase que el deudor tuviere a favor de su cónyuge, de sus parientes, o de
los gerentes, administradores, apoderados con poder general de administración u otras
personas con injerencia en la administración de sus negocios. La ley precisa que se entiende
por parientes a los ascendientes, descendientes y colaterales hasta el cuarto grado de
consanguinidad y afinidad, inclusive.

Quien planifica una reorganización debe asumir, en consecuencia, que el pasivo laboral
preferente seguirá litigándose durante la protección. Véase el capítulo
\ref{cap:creditos-laborales}.
\end{alerta}

\subsubsection{El estatuto laboral separado del artículo 60 A}

La \Lreforma{} \citeyearpar{ley-21563} incorporó el \art{60 A}, que zanja de modo expreso la
relación entre reorganización y derecho del trabajo: los derechos de los trabajadores con
contrato vigente, y los de aquellos cuyo contrato terminó manteniendo la empresa
obligaciones laborales y previsionales pendientes, \destacar{se rigen por las normas del
\cdt{} y demás normas que correspondan, sin que sean aplicables las de la ley concursal},
salvo lo dispuesto en la letra a) del \art{57}.

\begin{foco}[Consecuencia práctica del artículo 60 A]
La reorganización no reestructura el pasivo laboral. Los trabajadores no son acreedores
sujetos al acuerdo en las condiciones que éste fije: conservan su estatuto propio. Una
propuesta que pretenda imponerles quitas o esperas carece de sustento legal, y su sola
inclusión debilita la credibilidad del plan.
\end{foco}

\subsection{Duración y prórroga de la Protección Financiera Concursal}

El \art{57} \Nro 1 fija el plazo base en \destacar{sesenta días} contados desde la
notificación de la Resolución de Reorganización. La \Lreforma{} duplicó el plazo original de
treinta días.

El \art{58} regula la prórroga mediante tres vías alternativas, cuyos quórums difieren:

\begin{enumerate}
  \item \textbf{Primera prórroga:} hasta por sesenta días adicionales, con el apoyo de dos o
  más acreedores que representen más del 30\,\% del total del pasivo, excluidos los créditos
  de las personas relacionadas.

  \item \textbf{Segunda prórroga:} hasta por otros sesenta días, solicitada a más tardar el
  décimo día anterior al vencimiento del plazo vigente, con el apoyo de dos o más acreedores
  que representen más del 50\,\% del pasivo, excluidas las personas relacionadas.

  \item \textbf{Prórroga única:} en un solo acto, hasta por ciento veinte días, con el apoyo
  de dos o más acreedores que representen más del 50\,\% del pasivo, excluidas las personas
  relacionadas.
\end{enumerate}

El techo es, por cualquiera de los dos caminos, de \destacar{ciento ochenta días}.

\begin{practica}[El plazo de diez días que se pasa por alto]
La segunda prórroga debe solicitarse \emph{hasta el décimo día anterior} al vencimiento del
plazo en curso. No corre desde el vencimiento: corre hacia él. Quien espera a los últimos
días pierde la posibilidad de pedirla.
\end{practica}

\begin{foco}[Los acreedores con garantía real no pierden su preferencia]
El \art{58} precisa que los acreedores hipotecarios y prendarios que apoyen la prórroga no
pierden su preferencia y conservan la facultad de impetrar medidas conservativas. Es un
argumento decisivo al negociar con ellos: apoyar la prórroga no les cuesta su posición.
\end{foco}

\subsection{Verificación y objeción de créditos}

El \art{70} concede a los acreedores un plazo de \destacar{quince días} contado desde la
notificación de la Resolución de Reorganización para verificar sus créditos ante el
tribunal. La \Lreforma{} amplió el plazo original de ocho días.

Al verificar deben acompañarse los títulos justificativos, señalando si el crédito se
encuentra garantizado con prenda o hipoteca y el avalúo comercial de los bienes sobre los
que recaen las garantías.

\begin{practica}[Cuándo no es necesario verificar]
El propio \art{70} exime de verificación cuando el crédito y el avalúo comercial de las
garantías figuran, \emph{a satisfacción del acreedor}, en el estado de deudas acompañado por
el deudor. La expresión es relevante: la conformidad del acreedor es la que exime. Ante
cualquier discrepancia de monto, preferencia o avalúo, conviene verificar igualmente.
\end{practica}

\subsection{Contenido de la propuesta}

\subsubsection{Libertad de objeto}

El \art{60} establece que la propuesta puede versar sobre cualquier objeto tendiente a
reestructurar los pasivos y activos de la empresa. La amplitud es deliberada: quitas,
esperas, capitalización de créditos, enajenación de unidades productivas, constitución de
garantías, cambio de administración o combinaciones de todo ello.

\subsubsection{Clases o categorías de acreedores}

El \art{61} permite separar la propuesta en clases o categorías, formulando una para los
acreedores valistas y otra para los hipotecarios y prendarios cuyos créditos estén
garantizados con bienes del deudor o de terceros. Dos reglas la gobiernan:

\begin{enumerate}
  \item Los acreedores hipotecarios y prendarios que voten la propuesta \destacar{conservan
  sus preferencias}.
  \item La propuesta debe ser \destacar{igualitaria para todos los acreedores de una misma
  clase}, salvo acuerdo en contrario adoptado conforme al \art{64}.
\end{enumerate}

\begin{foco}[Por qué la división en clases suele ser la decisión que salva el acuerdo]
Los acreedores con garantía real y los valistas tienen expectativas de recuperación
radicalmente distintas en un escenario de liquidación. Someterlos a una propuesta única
obliga a ofrecer al valista más de lo razonable, o al garantizado menos de lo que aceptará.
La separación en clases permite calibrar cada oferta contra el escenario de liquidación
específico de esa clase, que es el único término de comparación relevante para cada grupo.
\end{foco}

\subsubsection{Posposición de las personas relacionadas}

El \art{63} pospone en el pago a los acreedores personas relacionadas con el deudor cuyos
créditos no se encuentren debidamente documentados \destacar{noventa días antes} del inicio
del procedimiento, hasta que se paguen íntegramente los créditos de los demás acreedores
afectados por el acuerdo. El acuerdo puede, además, extender la posposición a otras personas
relacionadas cuyos créditos sí estén documentados, previo informe fundado.

\begin{alerta}[Una fecha que conviene anotar]
El corte de los noventa días opera hacia atrás desde el inicio del procedimiento. Documentar
un crédito con persona relacionada en vísperas de presentar no evita la posposición, y sí
puede leerse como un acto dirigido a alterar el mapa de mayorías.
\end{alerta}

\subsubsection{Acreedores comprendidos en el acuerdo}

Según el \art{66}, los acuerdos sólo afectan a los acreedores cuyos créditos se originen con
anterioridad a la Resolución de Reorganización. Los créditos posteriores quedan fuera.

Los acreedores anteriores que no verificaron oportunamente, y los no contenidos en el
certificado del \art{55}, pueden demandar el cumplimiento del acuerdo a su favor mientras no
prescriban las acciones que de él resulten.

\begin{foco}[La regla que hay que explicar al cliente antes de presentar]
El acuerdo obliga a acreedores que no comparecieron y que el deudor pudo haber olvidado. Un
pasivo omitido no desaparece: reaparece exigiendo el cumplimiento del acuerdo a su favor,
con el consiguiente desajuste del flujo proyectado. Es la razón técnica por la que la nómina
de acreedores debe construirse con el informe consolidado y no con la memoria del cliente
(capítulo \ref{cap:entrevista}).
\end{foco}

\subsection{Financiamiento y disposición de activos durante la protección}

El \art{74}, en su redacción vigente tras la \Lreforma{} \citeyearpar{ley-21563}, habilita al
deudor, durante la protección, para:

\begin{itemize}
  \item \textbf{Enajenar activos} cuyo valor no exceda el \destacar{20\,\% de su activo fijo
  contable}.
  \item \textbf{Contratar préstamos u otras operaciones de financiamiento} para sus operaciones,
  siempre que no superen el \destacar{20\,\% de su pasivo} según la certificación contable del
  \art{55}, circunstancia que el veedor debe certificar.
\end{itemize}

Las operaciones que excedan esos límites, así como \destacar{toda operación con personas
relacionadas} cualquiera sea su monto, requieren la autorización de acreedores que representen
\destacar{más del 30\,\% del pasivo, excluidos los créditos de las personas relacionadas}. La
reforma rebajó aquí el umbral ---antes del 50\,\%--- y depuró la base de cálculo, facilitando el
acceso a financiamiento sin dejarlo en manos de los acreedores vinculados al deudor.

\begin{alerta}[El umbral del 20\,\% se calcula sobre bases distintas]
Para la venta de activos, el 20\,\% se mide sobre el \emph{activo fijo contable}. Para los
préstamos, sobre el \emph{pasivo} certificado. Confundir las bases produce operaciones que
exceden el límite sin autorización, con el riesgo de su posterior cuestionamiento.
\end{alerta}

\subsubsection{La preferencia del nuevo financiamiento y del suministro (\emph{new money})}

El incentivo que hace operativo el \art{74} está en su efecto ante un eventual fracaso: los
préstamos y operaciones de financiamiento contratados durante la protección \destacar{no se
computan en las nóminas de créditos} y, si el acuerdo no se suscribe y sobreviene la liquidación,
se pagan con la \destacar{preferencia de primera clase del \artcc{2472} \Nro 4}. Es la traducción
chilena del \emph{new money} o financiamiento post-solicitud: quien aporta liquidez a la empresa en
crisis obtiene un rango preferente que compensa el riesgo, y por eso el \art{63} excluye estos
créditos de la posposición de las personas relacionadas.

El \art{72} extiende una lógica gemela a la \destacar{continuidad del suministro}: los proveedores
de bienes y servicios necesarios para el funcionamiento de la empresa ---cuyos créditos anteriores
a la Resolución de Reorganización no superen en conjunto el 20\,\% del pasivo--- se pagan
\destacar{preferentemente en las fechas convenidas} siempre que \emph{mantengan el suministro} en
las mismas condiciones, lo que el veedor debe acreditar; una vez pagados, salen del pasivo con
derecho a voto, y si sobreviene la liquidación gozan también de la preferencia del \artcc{2472}
\Nro 4. La ley premia, así, al proveedor que no abandona a la empresa en reorganización.

\begin{practica}[El nuevo financiamiento y el proveedor fiel: dos preferencias que se negocian]
El financiamiento del \art{74} es, muchas veces, la única vía de capital de trabajo cuando el
crédito ordinario está cerrado; su exploración debe iniciarse el primer día de la protección, no
cuando la caja ya se agotó. Y el \art{72} ofrece una carta de negociación con los proveedores
críticos: mantener el suministro a cambio del pago preferente en los términos originales. Ambas
preferencias ---new money y suministro--- son herramientas activas del plan, no tecnicismos.
\end{practica}

\subsection{La junta llamada a pronunciarse sobre la propuesta}

\subsubsection{Quórum de aprobación}

El \art{79} fija la regla central del procedimiento. Cada clase o categoría se analiza,
delibera y acuerda \emph{en forma separada en la misma junta}. La propuesta se entiende
acordada cuando cuenta con:

\begin{enumerate}
  \item el \destacar{consentimiento del deudor}; y
  \item el voto conforme de \destacar{dos tercios o más de los acreedores presentes}, que
  representen \destacar{al menos dos tercios del total del pasivo con derecho a voto} de la
  respectiva clase o categoría.
\end{enumerate}

No pueden votar las personas relacionadas con el deudor.

\begin{foco}[La doble mayoría y su lectura estratégica]
El quórum exige simultáneamente mayoría de personas y mayoría de monto, ambas de dos
tercios, dentro de cada clase. La consecuencia es que ni el acreedor mayoritario aislado
puede imponer el acuerdo, ni una mayoría de acreedores pequeños puede hacerlo contra el
titular del grueso del pasivo. La negociación debe cubrir ambas dimensiones: número y monto.
\end{foco}

\subsubsection{Suspensión de la junta}

El \art{82} permite a la junta acordar, con quórum calificado, la suspensión por no más de
diez días, fijando nuevo día y hora para su reanudación. El deudor \destacar{conserva la
protección financiera} hasta la celebración de esa junta.

\begin{practica}[Un recurso de negociación de último minuto]
Cuando el recuento anticipa un rechazo por margen estrecho, la suspensión ofrece diez días
para cerrar la diferencia sin perder la protección. Es preferible suspender que forzar una
votación perdida, porque el rechazo no admite reversa.
\end{practica}

\subsubsection{Modificación del acuerdo}

Conforme al \art{83}, las modificaciones deben adoptarse por el deudor y los acreedores que
lo suscribieron, agrupados en sus clases, con el mismo procedimiento y mayorías del \art{79}.

\subsection{Impugnación del acuerdo}

\subsubsection{Causales}

El \art{85} enumera las causales taxativas por las que los acreedores afectados pueden
impugnar el acuerdo. Entre ellas:

\begin{enumerate}
  \item Defectos en las formas establecidas para la convocatoria y celebración de la junta,
  que hubieren impedido el ejercicio de los derechos de los acreedores o del deudor.
  \item Error en el cómputo de las mayorías, \destacar{siempre que incida sustancialmente en
  el quórum} del acuerdo.
  \item Falsedad o exageración del crédito, o incapacidad o falta de personería para votar
  de alguno de los acreedores que concurrieron a formar el quórum necesario ---si, excluido ese
  acreedor o la parte falsa o exagerada de su crédito, no se alcanza el quórum---.
  \item \destacar{Acuerdo entre uno o más acreedores y el deudor} para votar a favor, abstenerse
  o rechazar, con el fin de obtener una ventaja indebida respecto de los demás acreedores.
  \item Ocultación o exageración del activo o del pasivo.
  \item Contener \destacar{una o más estipulaciones contrarias a la ley} o al ordenamiento
  jurídico.
\end{enumerate}

\begin{foco}[El estándar de sustancialidad]
El numeral 2 no se satisface con acreditar un error de cómputo: exige que el error
\emph{incida sustancialmente} en el quórum. Un error que no altera el resultado no invalida
el acuerdo. La impugnación debe, por tanto, demostrar el recálculo y su efecto sobre la
aprobación.
\end{foco}

\subsubsection{Plazo}

El \art{86} concede \destacar{cinco días} contados desde la publicación del acuerdo en el
\boletin{}. Las impugnaciones presentadas fuera de plazo \destacar{se rechazan de plano}.

\begin{alerta}[Cinco días, sin excepciones]
Es uno de los plazos más breves de la ley y su incumplimiento no admite subsanación. La
publicación en el \boletin{} es la única señal: no hay notificación personal. Quien
representa a un acreedor disidente debe tener el control diario del Boletín activado desde
la fecha de la junta.
\end{alerta}

\subsubsection{Audiencia única}

El \art{87} dispone que las impugnaciones se tramitan como un solo incidente y se fallan
conjuntamente en una audiencia única, que el tribunal cita dentro de los diez días de
vencido el plazo para impugnar. La audiencia es verbal y se celebra con los que asistan; en
ella deben resolverse las incidencias que promuevan las partes. El tribunal puede
suspenderla y continuarla con posterioridad. La resolución debe dictarse a más tardar dentro
de los treinta días siguientes a la fecha de celebración.

\subsubsection{Aprobación y vigencia}

Según el \art{89}, el acuerdo se entiende aprobado y comienza a regir una vez vencido el
plazo para impugnarlo sin que se haya impugnado, previa declaración del tribunal de oficio o
a petición de cualquier interesado o del veedor. Si fue impugnado y las impugnaciones se
desecharon, el tribunal lo declara aprobado en la misma resolución, y el acuerdo rige desde
que ésta causa ejecutoria. Ambas resoluciones se notifican en el \boletin{}.

\subsection{Cláusulas especiales del acuerdo}

Además de la libertad de objeto del \art{60}, la ley regula estipulaciones específicas que
conviene conocer al diseñar la propuesta.

\subsubsection{Diferencias entre acreedores de igual clase}

El \art{64} permite establecer condiciones \destacar{más favorables para algunos acreedores de
una misma clase}, siempre que los demás acreedores de esa clase lo acuerden con quórum
especial, calculado únicamente sobre el monto de los créditos de estos últimos. Es la
excepción a la regla de igualdad intraclase del \art{61}.

\subsubsection{Constitución de garantías y valorización de activos}

El \art{65} habilita a estipular la constitución de garantías dentro del acuerdo. Y el
\art{76} entrega al veedor la \destacar{valorización de los activos} a enajenar durante la
ejecución, así como la fiscalización de que el producto de esas operaciones ingrese
efectivamente a la caja de la empresa y se destine sólo a financiar su giro.

\subsubsection{Cláusula arbitral}

El \art{68} permite incorporar, en cualquier clase o categoría, una \destacar{cláusula
arbitral}: las diferencias sobre aplicación, interpretación, cumplimiento, terminación o
declaración de incumplimiento del acuerdo se someten a arbitraje (véase el capítulo
\ref{cap:arbitraje}).

\subsection{Efectos del acuerdo aprobado}

\subsubsection{Mérito ejecutivo}

El \art{90} confiere \destacar{mérito ejecutivo} al acuerdo: una copia del acta de la junta con
el voto favorable y el texto íntegro, junto con la resolución que lo aprueba y su certificado
de ejecutoria, autorizada ante ministro de fe o protocolizada, vale como título ejecutivo para
todos los efectos legales.

\subsubsection{Alcance subjetivo: obliga a todos}

El \art{91} fija el efecto vinculante: el acuerdo aprobado \destacar{obliga al deudor y a todos
los acreedores de cada clase, hayan o no concurrido} a la junta que lo acordó. Es la
manifestación de la fuerza colectiva del concurso: el disidente y el ausente quedan igualmente
sujetos.

\subsubsection{Efecto sobre los créditos y su tratamiento tributario}

El \art{93} dispone que los créditos comprendidos en el acuerdo se entienden \destacar{remitidos,
novados o repactados}, según corresponda, para todos los efectos legales. La norma agrega una
consecuencia tributaria de primer orden: el acreedor contribuyente de primera categoría puede
\destacar{deducir como gasto necesario} ---conforme al artículo 31 \Nro 4 de la Ley sobre Impuesto a la Renta\index[leyes]{DL 824 (Ley sobre Impuesto a la Renta)!art. 31 N.º 4}--- las cantidades
correspondientes a la condonación o remisión pactada en el acuerdo.

\begin{foco}[El incentivo tributario del artículo 93 facilita el acuerdo]
Que la quita sea deducible como gasto cambia el cálculo del acreedor financiero: la condonación
no es pérdida pura, sino gasto que reduce su base imponible. Este dato ---frecuentemente omitido
en la negociación--- puede ser el argumento que incline a un banco a aprobar una propuesta con
quita. Debe estar sobre la mesa al negociar.
\end{foco}

\subsection{Bienes no esenciales para la continuidad del giro}

El \art{94} permite al acreedor con prenda o hipoteca solicitar, dentro de los \destacar{quince
días} siguientes a la publicación de la Resolución de Reorganización, que el tribunal declare
que el bien garantizado \destacar{no es esencial} para el giro ---con la misma lógica del
\art{286 Q} del procedimiento simplificado (capítulo \ref{cap:reorganizacion-simplificada})---.
La calificación determina si el acreedor garantizado queda o no sujeto a los términos del
acuerdo respecto de ese bien.

\subsection{Rechazo del acuerdo: la liquidación de oficio}

El \art{96} regula la consecuencia del fracaso: si la propuesta es rechazada por falta de
quórum o porque el deudor no otorga su consentimiento, el tribunal dicta la
\destacar{Resolución de Liquidación de oficio y sin más trámite, en la misma junta}, salvo que
ésta disponga lo contrario por quórum especial. En ese caso, el deudor debe publicar, a través
del veedor, una nueva propuesta.

\begin{alerta}[La reorganización ordinaria sí tolera un segundo intento, a diferencia de la simplificada]
Aquí sí existe una válvula: la junta puede, por quórum especial, evitar la liquidación inmediata
y abrir espacio a una nueva propuesta. Es una diferencia importante con la reorganización
simplificada, donde el rechazo conduce a la liquidación sin esa alternativa (capítulo
\ref{cap:reorganizacion-simplificada}, \art{286 S}). Preparar el quórum especial para sostener
un segundo intento es, por tanto, una estrategia disponible en el procedimiento ordinario.
\end{alerta}

\subsection{Nulidad del acuerdo: una acción distinta de la impugnación}

El \art{97} regula la nulidad del acuerdo, que no debe confundirse con la impugnación del
\art{85}. Mientras la impugnación se funda en causales formales y se deduce dentro de los cinco
días de la publicación, la \destacar{nulidad} sólo procede por \destacar{ocultación o
exageración del activo o del pasivo} de las que se haya tomado conocimiento \emph{después} de
vencido el plazo para impugnar.

Tres notas: la declaración de nulidad \destacar{extingue de pleno derecho las cauciones} que
garantizan el acuerdo; la acción puede interponerla cualquier interesado; y \destacar{prescribe
en un año} contado desde que el acuerdo comenzó a regir.

\begin{foco}[Impugnación y nulidad: dos ventanas temporales distintas]
La impugnación ataca vicios visibles en los cinco días siguientes a la publicación. La nulidad
ataca el fraude \emph{oculto} que emerge después, con un año de plazo. Para el acreedor que
descubre tardíamente que el deudor exageró su pasivo para inflar mayorías o escondió activos, la
nulidad ---no la impugnación, ya precluida--- es la vía.
\end{foco}

\subsection{El Acuerdo de Reorganización Extrajudicial o Simplificado}

Junto a la reorganización judicial, el Título 3 del Capítulo III regula el \destacar{Acuerdo de
Reorganización Extrajudicial o Simplificado} (ARES): un acuerdo que la empresa deudora negocia
\emph{fuera} del tribunal con sus acreedores y luego somete a aprobación judicial. No debe
confundirse con la reorganización simplificada de MIPEs del capítulo
\ref{cap:reorganizacion-simplificada}, que es un procedimiento judicial distinto.

\subsubsection{Legitimación, competencia y quórum}

El \art{102} legitima a toda empresa deudora para celebrar el ARES. El \art{103} atribuye la
competencia al tribunal que habría conocido de una reorganización judicial. La clave está en el
\art{109}: el deudor debe presentar el acuerdo \destacar{suscrito por dos o más acreedores que
representen al menos las tres cuartas partes del pasivo} de la respectiva clase. Las personas
relacionadas \destacar{no pueden suscribirlo ni se cuentan} para el quórum.

\subsubsection{Normas aplicables y aprobación}

El \art{106} hace aplicables al ARES ---en lo que no contravenga su regulación propia--- las
normas de los acuerdos por clases, determinación del pasivo, condonación, garantías, cláusula
arbitral e interventor de la reorganización judicial. El \art{112} regula la aprobación: dentro
de los diez días de publicado el acuerdo, el tribunal puede citar a los acreedores afectados
para su aceptación con el quórum del \art{109}; aceptado, o vencido el plazo, se aprueba.

\begin{foco}[El ARES premia el acuerdo previo con un quórum más alto]
El ARES invierte el orden habitual: primero se negocia y se firma con las tres cuartas partes
del pasivo, después se pide la aprobación judicial. El quórum es más exigente que el de la
reorganización judicial (dos tercios), pero a cambio el procedimiento es más rápido y menos
expuesto: quien llega con el acuerdo ya cerrado evita la incertidumbre de la junta. Es la vía de
elección cuando el deudor tiene pocos acreedores financieros y relación fluida con ellos.
\end{foco}

\subsection{El incumplimiento del acuerdo y sus consecuencias}

Aprobado y vigente el acuerdo, resta la fase de cumplimiento. El \art{98} regula la
\destacar{acción de incumplimiento}: cualquier acreedor afectado puede solicitar que el acuerdo
se declare incumplido por inobservancia de sus estipulaciones, o \destacar{si se agrava el mal
estado de los negocios} del deudor de forma que haga temer un perjuicio para los acreedores.

El \art{99} sujeta las acciones de nulidad e incumplimiento al \destacar{juicio sumario} ante el
mismo tribunal que tramitó el acuerdo. La resolución que las acoge es apelable \destacar{en ambos
efectos} ---excepción al régimen recursivo restrictivo---, pero el deudor queda de inmediato
sujeto a la intervención de un veedor con facultades de interventor.

La consecuencia es drástica: conforme al \art{100}, firme la declaración de nulidad o
incumplimiento, el mismo tribunal dicta la \destacar{Resolución de Liquidación de oficio y sin
más trámite}. El \art{101} dispone que el liquidador es el propuesto en la demanda respectiva.

\begin{foco}[El incumplimiento del acuerdo conduce directo a la liquidación]
El acuerdo aprobado no es una salvación definitiva: su incumplimiento reabre el camino a la
liquidación, ahora sin juicio de oposición. Para el acreedor, la acción del \art{98} es la
garantía de que la quita o espera pactada se respete; para el deudor, es la advertencia de que
el acuerdo debe cumplirse al pie, porque el margen de tolerancia es escaso y la sanción, la
liquidación inmediata.
\end{foco}

\subsection{El papel del veedor en la negociación}

Conviene subrayar la función que el \art{25} asigna al veedor: su tarea principal es
\destacar{propiciar los acuerdos} entre el deudor y sus acreedores, facilitando la proposición y
negociación. Puede citar a deudor y acreedores en cualquier momento, desde la Resolución de
Reorganización hasta la entrega de su informe. El veedor no es un mero fiscalizador: es un
facilitador de la negociación, y usar ese rol ---convocar, mediar, acercar posiciones--- es parte
de la estrategia de quien busca aprobar el acuerdo.

\subsection{Espejo comparado: la reestructuración preventiva europea}

La reorganización chilena admite comparación con el modelo que la \destacar{Directiva (UE)
2019/1023} \parencite{ue-directiva-2019-1023} impuso a los Estados miembros y que España
transpuso en la \leyn{16/2022}{reforma del texto refundido concursal} \parencite{esp-ley16-2022}.

Dos rasgos del modelo europeo interesan al derecho chileno. Primero, la
\destacar{reestructuración preventiva}: los planes pueden adoptarse \emph{antes} de la
insolvencia declarada, cuando la crisis es aún reversible ---un umbral más temprano que el de la
reorganización chilena, que exige ya la calidad de empresa deudora en dificultades---. Segundo, y
más relevante, el \destacar{arrastre de clases disidentes} (\emph{cross-class cram-down}): el plan
puede imponerse a clases enteras de acreedores que lo rechazaron, si se cumplen las garantías de
la directiva ---prueba del interés superior de los acreedores y respeto de la prelación---.

\begin{foco}[El arrastre de clases: lo que Chile no tiene]
El sistema chileno exige el consentimiento de cada clase por su doble quórum de dos tercios
(\art{79}): una clase que rechaza bloquea el acuerdo respecto de ella. El modelo europeo permite
\emph{arrastrar} a la clase disidente si el plan es equitativo y le ofrece más que la
liquidación. Es la diferencia estructural: donde Chile requiere acuerdo de todas las clases,
Europa admite imponer el plan a la que se opone. La ausencia de arrastre en Chile da a cada
clase un poder de veto que condiciona toda la negociación.
\end{foco}

%==============================================================================
\bloquePractico
%==============================================================================

\subsection{Iter procesal de la reorganización judicial ordinaria}

\begin{flujograma}{Tramitación de la reorganización judicial ordinaria}{cap04-iter}
  \node[hito] (ingreso) {INGRESO DE SOLICITUD DE REORGANIZACIÓN\\[2pt]
    \normalfont Modelo de la NCG N.\textsuperscript{o} 4 (\art{54})};
  \node[nodo, below=of ingreso] (resol)
    {RESOLUCIÓN DE REORGANIZACIÓN (\art{57})\\[2pt]
     \normalfont Dentro del quinto día. Inicio de la PFC: 60 días};
  \node[nodo, below=of resol] (verif)
    {VERIFICACIÓN DE CRÉDITOS (\art{70})\\[2pt]
     \normalfont Plazo de 15 días desde la notificación};
  \node[nodo, below=of verif] (informe)
    {INFORME DEL VEEDOR SOBRE LA PROPUESTA};
  \node[hito, below=of informe] (junta)
    {JUNTA DE ACREEDORES (\art{79})\\[2pt]
     \normalfont Doble quórum de dos tercios, por clase};

  \coordinate (split) at ($(junta.south)+(0,-6mm)$);
  \node[favorable, below left=12mm and 4mm of junta.south] (aprueba)
    {ACUERDO ADOPTADO};
  \node[adverso, below right=12mm and 4mm of junta.south] (rechaza)
    {RECHAZO DE LA PROPUESTA};
  \node[nodo, text width=40mm, below=9mm of aprueba] (impug)
    {Impugnación: 5 días desde la publicación (\art{86})};
  \node[adverso, text width=40mm, below=9mm of rechaza] (liq)
    {APERTURA DE LA LIQUIDACIÓN};

  \draw[flecha] (ingreso) -- (resol);
  \draw[flecha] (resol) -- (verif);
  \draw[flecha] (verif) -- (informe);
  \draw[flecha] (informe) -- (junta);
  \draw[flecha] (junta.south) -- (split) -| (aprueba.north);
  \draw[flecha] (junta.south) -- (split) -| (rechaza.north);
  \draw[flecha] (aprueba) -- (impug);
  \draw[flecha] (rechaza) -- (liq);
\end{flujograma}

\subsection{Calendario de plazos del procedimiento}

\begin{table}[htbp]
\centering
\small
\caption{Plazos de la reorganización judicial ordinaria}
\label{tab:plazos-reorganizacion}
\begin{tabularx}{\textwidth}{@{}>{\raggedright\arraybackslash}p{0.34\textwidth} p{0.17\textwidth} X@{}}
\toprule
\textbf{Actuación} & \textbf{Plazo} & \textbf{Norma} \\
\midrule
Dictación de la Resolución de Reorganización & 5.\textsuperscript{o} día desde la
presentación & \art{57} \\
\addlinespace
Protección Financiera Concursal (base) & 60 días & \art{57} \Nro 1 \\
\addlinespace
Primera prórroga (apoyo $>$30\,\% del pasivo) & Hasta 60 días & \art{58} \\
\addlinespace
Segunda prórroga (apoyo $>$50\,\%) & Hasta 60 días & \art{58} \\
\addlinespace
Prórroga única (apoyo $>$50\,\%) & Hasta 120 días & \art{58} \\
\addlinespace
Solicitud de la segunda prórroga & Hasta el 10.\textsuperscript{o} día \emph{anterior} al
vencimiento & \art{58} \\
\addlinespace
Verificación de créditos & 15 días desde la notificación & \art{70} \\
\addlinespace
Suspensión de la junta & Hasta 10 días & \art{82} \\
\addlinespace
Impugnación del acuerdo & 5 días desde la publicación & \art{86} \\
\addlinespace
Citación a la audiencia de impugnaciones & 10 días de vencido el plazo para impugnar
& \art{87} \\
\addlinespace
Resolución de las impugnaciones & 30 días desde la audiencia & \art{87} \\
\bottomrule
\end{tabularx}
\end{table}

\subsection{Construcción del mapa de mayorías}

Dado el doble quórum de dos tercios del \art{79}, el trabajo previo consiste en construir,
por cada clase, una matriz de dos columnas: número de acreedores y monto del pasivo con
derecho a voto. Sobre ella se identifican:

\begin{checklist}[Elementos de la matriz]
  \item Acreedores que deben excluirse por ser personas relacionadas.
  \item Créditos sujetos a posposición conforme al \art{63}.
  \item Acreedores con garantía real y su clase separada.
  \item Umbral de dos tercios en número y en monto, calculado por clase.
  \item Acreedores cuyo voto es indispensable para alcanzar el umbral de monto.
  \item Acreedores necesarios para alcanzar el umbral de número.
  \item Estimación de la recuperación de cada clase en escenario de liquidación, como
  término de comparación de la oferta.
\end{checklist}

\subsection{Actuaciones del veedor y relación de trabajo}

El veedor no representa al deudor ni a los acreedores: informa. Las instrucciones que rigen
su actuación en los procedimientos de reorganización están contenidas en el Instructivo
SUPERIR \Nro 3 \parencite{instructivo-i3}, cuya lectura permite anticipar qué información
requerirá y en qué oportunidad.

\begin{practica}[Anticipación como estrategia]
Entregar al veedor la información antes de que la solicite, de forma ordenada y completa,
produce un informe más favorable que responder de modo reactivo a sus requerimientos. El
informe influye en el voto, y el voto decide el procedimiento.
\end{practica}

%==============================================================================
\bloqueJurisprudencial
%==============================================================================

\subsection{Límites del control judicial de mérito del acuerdo}

El acuerdo aprobado por las mayorías legales expresa una decisión económica de los
acreedores. La cuestión es si el tribunal puede revisar el mérito de esa decisión o si su
control se limita a la legalidad del procedimiento y al cumplimiento de los quórums.

El sistema de causales taxativas del \art{85} sugiere la segunda respuesta: las causales de
impugnación son formales o de cómputo, no de contenido. \pendiente{Sistematizar la
jurisprudencia de la Sala Civil de la \csup{} sobre el alcance del control judicial.}

\subsection{Quitas de capital y acreedores disidentes con garantía real}

Una de las cuestiones dogmáticas más delicadas de la reorganización es hasta dónde puede la
mayoría imponer una \destacar{quita} ---remisión parcial del capital--- al acreedor
\destacar{disidente con garantía real}. El punto de partida es el \art{61}: la propuesta puede
separarse en clases, y los acreedores hipotecarios y prendarios que \destacar{voten} su propuesta
\emph{conservan sus preferencias}. La conservación de la preferencia está, así, condicionada a la
concurrencia del acreedor garantizado a la formación del acuerdo de su clase.

El problema surge con el garantizado que \emph{no} vota o que disiente. Aquí se enfrentan dos
lecturas. La primera, \destacar{contractualista}, sostiene que el acuerdo es una convención cuyo
efecto vinculante sobre el disidente ---la fuerza de la mayoría--- no puede extenderse a sacrificar
una garantía real constituida antes del concurso sin vaciar de contenido el derecho de prenda
general reforzado; la quita sobre el crédito garantizado disidente sería, en esta tesis, ilegítima
más allá del valor de realización de la garantía. La segunda, \destacar{concursalista}, subraya que
la separación por clases del \art{61} y la regla de trato igualitario intraclase del \art{64}
suponen precisamente que la clase de garantizados delibere y quede vinculada por su mayoría interna:
lo que la ley protege es la \emph{preferencia} ---el derecho a pagarse con el producto del bien---,
no la intangibilidad nominal del crédito.

La lectura más equilibrada, coherente con el derecho comparado (el \emph{cross-class cram-down} de
la Directiva \textcite{ue-directiva-2019-1023} y la \emph{absolute priority rule} estadounidense),
distingue: la mayoría de la clase de garantizados puede imponer una quita al disidente de su clase,
pero \destacar{con el límite del valor de la garantía} ---nadie puede ser forzado a recibir menos de
lo que obtendría en la liquidación de su colateral (test del \emph{best interest of creditors})---.
Por debajo de ese piso, la quita degrada la preferencia y es impugnable; por encima, expresa la
decisión legítima de la clase. El \art{61} chileno, al anudar conservación de preferencia y
votación, se deja leer en esta clave.

\begin{foco}[El piso de la garantía como límite de la quita]
La discusión no es si el garantizado disidente puede sufrir una quita, sino cuánto. El umbral
defendible es el \destacar{valor de realización de su garantía}: hasta ahí su preferencia es
intangible; el excedente de su crédito, no cubierto por el bien, es materia valista y sigue la
suerte de esa clase. Al negociar un acuerdo con garantizados reticentes, valorizar bien el
colateral no es un detalle técnico: fija el terreno de lo que la mayoría puede o no imponer.
\end{foco}

\subsection{Impugnación por abuso del derecho de acreedores relacionados}

\pendiente{Sistematizar los fallos sobre impugnación fundada en la participación de personas
relacionadas en la formación del quórum, en relación con la prohibición de voto del \art{79}
y la posposición del \art{63}.}
